\chapter{The Self-Bootstrap and the First Theorem}
\label{chap:self-bootstrap-first-theorem}

The first seven chapters built an instrument, piece by piece: a tower of characteristics --- being
distinguishable, being admissible, being countable, and all the rest --- each defined on top of the
ones below it. This chapter does three things to finish the instrument and hand it over to the rest of
the book. It shows that the tower is not empty, by having the construction build something that
satisfies every characteristic at once. It extends the apparatus with the higher notions a working
science needs --- knowledge, reality, locality, calibration, measurement. And it turns the finished
apparatus on the one thing it was built to study, truth itself, and proves its first theorem. That
last act is a hinge: everything before it builds the instrument, and everything after it uses the
instrument to do mathematics.

\section{The tower inhabits itself}
\label{se:the-tower-inhabits-itself}

A definition is a promise that can be empty. To say what it would mean for an object to be
distinguishable, admissible, and countable all at once is not yet to exhibit one, and a tower of
characteristics that nothing satisfies is an elaborate way of describing nothing. The construction
closes this gap directly. It builds an instance of every characteristic in the tower, in order, each
constructed from the instances below it, so that by the top of the chain there is a single object that
provably possesses the whole stack of properties. The tower is shown to be inhabited, and it is
inhabited by something the construction assembled out of its own earlier pieces. This is the
self-bootstrap, and it is best read as a funge. The earlier chapters tanged the tower apart, defining
each characteristic as a distinct selection; the bootstrap funges those rungs back into one. The
construction places a single seed at the bottom --- the trivial Fact --- and the property of satisfying
each characteristic propagates upward from it, pooling the many rungs into a single inhabited instance
under the one characteristic of being reachable from the seed. The apparatus does not merely describe a
kind of object; it produces one, built from the same Facts that defined the kind.

It is worth walking the bootstrap one rung at a time, because the order is the whole point. The
construction starts where the first chapter started, with the trivial Fact that true equals true, and
shows it is distinguishable --- there is a decidable way to tell it from something else. Having a
distinguishable object in hand, it shows that object is admissible, the property built directly on
distinguishability; having an admissible one, it shows it is countable; and so on up, each rung
reached from the witness the previous step produced, the single seed carried all the way up. No rung is
asserted; each is constructed from the one below. By the top of the chain the same object, elaborated
through every step, satisfies the entire tower at once, and the construction has its proof that the
tower describes something rather than nothing --- a proof that is itself an object, the inhabited tower
made concrete.

The maneuver has the flavor of a system certifying itself, and the construction is candid that this is
a delicate thing to do. The experiment it leans on here is the prover--verifier effect, which draws
the line carefully. An informational theory like this one, the experiment says, is not complete in
isolation; it functions as a \emph{verifier}, checking whether proposed structures are admissible,
and it requires the full range of admissible physical models to act as its \emph{prover}, supplying
the structures to be checked. The ledger of recorded Facts is the unique fixed point where prover and
verifier agree. Read against the self-bootstrap, this is exactly the right humility: the construction
can verify that an object satisfies its tower of characteristics, and it can build such an object, but
it does not claim to generate all truth from nothing. It certifies; the world proposes; the inhabited
tower is the fixed point where the two meet.

There is a worry the self-bootstrap has to answer, the worry that haunts every self-describing system:
has the construction smuggled in what it claims to build? The answer is in the strict order of the
chain. Each instance is constructed before the characteristic that consumes it is ever invoked, so at no
point does a step depend on a witness that has not already been exhibited. The construction never
assumes the tower is inhabited in order to prove it inhabited; it walks up from the trivial Fact, which
costs nothing and assumes nothing, and adds one exhibited witness at a time. That strictness is what
separates an honest bootstrap from a circular one, and it is the reason the inhabited tower counts as a
proof rather than a wish.

\section{Extending the apparatus}
\label{se:extending-the-apparatus}

With the tower inhabited, the construction adds the characteristics a science needs above bare
counting and comparison, and it adds them the same way it added everything else --- as decidable
properties layered on the ones below. It builds a notion of knowledge, of what it takes for a recorded
Fact to count as something known rather than merely registered. It builds reality and locality: a Fact
is treated as real when it has been witnessed, and as local when its witnessing does not depend on
events outside its own neighborhood, which is the fifth chapter's local present promoted to a property
an object can have. It builds paths through space and time and a heartbeat that steps them forward.
And it builds the apparatus of measurement proper --- distance, speed, and the property of having been
measured at all.

Two of these additions are worth a second look, because they make precise ideas the earlier chapters
only gestured at. Reality, here, is not a metaphysical status but a bookkeeping one: a Fact is real when
it has been witnessed --- entered into some carrier's record --- and a description that posits
unwitnessed facts, sub-resolution structure no instrument ever recorded, is declined exactly as the
first chapter's Berkeley--Galileo rule declined it. Locality is the fifth chapter's local present
hardened into a property an object carries: a Fact is local when deciding it needs only its own
neighborhood, with no reach into events a finite signal could not yet have touched. Between them,
reality and locality fence the apparatus off from two temptations at once --- inventing facts that were
never recorded, and letting a here depend on a there it cannot have reached --- and both fences are the
same discipline the construction has enforced since its first page: nothing counts until it has been
produced.

One addition deserves singling out, because it is where the instrument is tuned. The construction
gives itself a calibration: it runs a known reference trace through the apparatus and checks that the
apparatus reports it correctly, the way a clinician checks a monitor against a known signal before
trusting it on a patient. The experiment that names the discipline here is the one on precision, and
its central distinction is one every experimentalist lives by. \emph{Precision} is the determinacy of
the values the apparatus supplies: once the anchor points are fixed, the minimal-curvature completion
assigns a definite value at every point between them, and that definiteness is precision. \emph{Accuracy}
is something else entirely --- agreement with future measurements --- and no amount of precision
guarantees it. The construction can be perfectly precise about the completion it draws through its
anchors and still be wrong about the next event, and keeping those two notions apart is what calibration
is for. A calibrated instrument is one whose precision has been checked against a reference; whether it
is accurate is a question only the future can answer.

The distinction is worth a concrete case, because conflating the two is the commonest error in
measurement. Suppose an instrument is anchored on three recorded events and draws its minimal-curvature
completion through them. Ask it for the value halfway between the first two anchors and it returns a
single, definite number --- it is perfectly precise, and asking twice gives the same answer to the last
digit. Now let the world produce a fourth event at that halfway point. If the event matches the
completion, the instrument was accurate there; if it does not, the instrument was exactly as precise as
before and simply wrong. Precision was a property of the completion the instrument drew; accuracy was a
verdict the world delivered later. Calibration cannot manufacture accuracy --- nothing can, in advance
--- but it can certify that the precision is trustworthy, that the completion is the unique one the
anchors admit, so that when the instrument is wrong it is wrong honestly and not because it was sloppy.

\section{Pointing the apparatus at truth}
\label{se:pointing-the-apparatus-at-truth}

Now the instrument is complete, and the construction does the thing the whole book has been building
toward: it points the apparatus at the very notion it began with, the truth of a proposition, and
proves a theorem about it. It builds a carrier whose values are propositions themselves --- so that a
proposition can be measured the way a number was --- and on that carrier it establishes its first
genuine theorem rather than its first construction.

The theorem is stated provocatively, and it is worth reading carefully because the provocation hides a
precise and modest claim. The construction calls it, in effect, the theorem that true equals false ---
and it does not say what that phrase seems to say. It does not prove that the trivial truth is the same
as falsehood, which would collapse the whole edifice. What it proves is a statement about
proof-irrelevance. A proposition is a \emph{subsingleton} when it has at most one proof --- when, as
far as the bookkeeping is concerned, all its proofs are interchangeable and only its yes-or-no status
matters. The theorem observes that the truth of the trivial Fact is a subsingleton, that the
\emph{negation} of that truth is also a subsingleton, and that these two facts are, as propositions,
equal. Measured by the only thing this apparatus measures --- how many distinguishable proofs a claim
carries --- true and not-true come out the same: each carries no informative multiplicity, one
trivially inhabited and the other trivially empty, both flat. The first theorem is the discovery that
at the level of proof-counting, the difference between a truth and its denial is not itself a
measurable distinction.

It helps to see what the theorem funges together and what it leaves apart. It does not pool true and
false as \emph{values} --- a Fact that decides true and one that decides false are as different as they
ever were, and every earlier chapter's arithmetic depends on the difference. What it pools is something
thinner: the two propositions \emph{that the truth is proof-irrelevant} and \emph{that the non-truth is
proof-irrelevant}. Both are flatly true, neither carries any internal multiplicity to distinguish it,
and so as objects in the proof-bookkeeping they are one and the same. The theorem funges the two by
their shared characteristic --- carrying no informative proof-structure --- while leaving the underlying
truth-values exactly as separate as before. It is a statement about the grain of the apparatus: below a
certain level of proof-detail, the apparatus cannot tell a flat yes from a flat no, and the theorem
records precisely where that floor lies.

Two things about how it is proved matter as much as what it says. First, it is established by
proposition extensionality --- the principle that propositions with the same truth-content are the
same proposition --- which is one of the construction's few standing assumptions, named openly. Second,
and more telling, the proof is \emph{constructive}: in reaching it the construction deleted its one
remaining appeal to the law of the excluded middle, the assumption that every proposition is either
true or false with no third option. It had been using that law to inhabit a generic either-or, found
that no constructive witness existed for the general case and that nothing actually depended on it, and
removed it. This is the same discipline that let the construction do without a free-floating power to
choose: wherever a result can be reached by building a witness rather than by assuming one exists, the
construction insists on building it. The first theorem is small, but it is honest all the way down ---
proved about the apparatus's own foundations, with the excluded middle taken out by the root.

This is worth dwelling on, because it is the quiet thesis of the whole first part. A finite construction
is not entitled to the two classical conveniences that let ordinary mathematics move faster --- the
assumption that every proposition is decidably true or false, and the assumption that one can choose an
element from each of infinitely many collections without saying how. Both assert that something exists
without exhibiting it, and a construction that can only ever exhibit has no use for either. The first
theorem is the first place this shows up as a deliberate deletion rather than a mere abstinence: the
construction had reached for the excluded middle, caught itself, and found it could finish without.
Everything in the second part of the book --- the variational results, the field theory, the eventual
electron --- is built under the same rule, which is why, when those results come to lean on a single
genuine assumption later, that assumption will stand out so sharply against the constructive background.

\section{What is demonstrated, and the turn to Act II}
\label{se:demonstrated-assumed-ch8}

This chapter closes the first part of the book, so the separation of built from assumed doubles as a
boundary marker.

What is demonstrated is that the instrument is finished and works. The tower of characteristics is
inhabited by an object the construction built from its own pieces; the apparatus has been extended with
knowledge, reality, locality, paths, calibration, and measurement, each a decidable property; and the
apparatus, turned on the truth of a proposition, proves a first theorem --- that truth and its negation
are indistinguishable by proof-count --- constructively, with proposition extensionality named and the
excluded middle removed. Everything here is finite and checkable, and the constructive discipline is
the genuine achievement, not a technicality.

One change is worth marking, because the reader will feel it in the next chapter. Up to here the
construction has been writing itself --- inventing custom characteristics, stacking classes, naming each
new capability as it went. From the next chapter on its own writing grows simpler and more ordinary: the
elaborate tower of bespoke properties gives way to plain structures and integers, and its output turns
from apparatus into theorems proved outright. The reason is exactly the one this chapter earned. The
instrument is finished; there is no more instrument to build, only mathematics to do with it, and
mathematics done with a finished instrument looks like ordinary mathematics --- which is the strongest
evidence the instrument was worth building.

\begin{quote}
\emph{A note on the labels.} The higher characteristics of this chapter carry the heaviest everyday
words the book has used --- \emph{knowledge}, \emph{reality}, \emph{locality}, \emph{measurement} ---
and each is a definition the construction stipulates, not a claim it proves about the world. To call a
witnessed Fact \emph{real} or a checked completion \emph{calibrated} is to fix a meaning, useful
exactly insofar as it matches how scientists already use the words. The one result that is a genuine
theorem rather than a naming is the first theorem itself, and even there the provocative title is a
label over a modest, exactly stated fact about subsingletons.

\end{quote}

What is assumed, then, is the short list the construction has carried openly from the start, now joined
by proposition extensionality and unburdened of the excluded middle. With the instrument complete, the
character of the work changes, and the next chapter changes with it. From here the construction stops
building measuring apparatus and starts using it to do mathematics --- the source itself grows simpler,
trading its tower of custom characteristics for ordinary structures and integers and fully proved
theorems. The instrument is put down, and the thing it was built to find is picked up: the single
invariant that the whole middle of the book is named for, beginning with the question of which path a
measurement actually takes when it is free to vary.
